% ============================================================
%  Section 3 — A Model of Collector and Consumer Demand
% ============================================================

\section{A Model of Collector and Consumer Demand}
\label{sec:model}

We develop a parsimonious auction model with heterogeneous buyers to
rationalize the non-monotone price-age profile documented in
Section~\ref{sec:results}.  The model builds on the two-type framework
of \citet{lovo2018}, adapted to the distinctive feature of wine:
a biological maturity date after which consumption value declines.

%% ── 3.1  Environment ────────────────────────────────────────────

\subsection{Environment}
\label{sec:model:env}

Consider a single bottle of wine with vintage-specific maturity date
$\tau > 0$ (years since vintage).  The bottle reaches auction at age
$a \geq 0$; we work throughout with \emph{normalised age}
$a^{*} \equiv a / \tau$, so that $a^{*} = 1$ corresponds to peak
maturity.

In each auction, $n \geq 2$ risk-neutral potential bidders compete in a
second-price sealed-bid (Vickrey) mechanism.  Each bidder is
independently drawn as one of two types:
\begin{itemize}
  \item \textbf{Consumption buyers} (type~$c$), with probability
        $\lambda(a^{*})$, value wine for drinking.
  \item \textbf{Collector buyers} (type~$k$), with probability
        $1 - \lambda(a^{*})$, value wine for holding, prestige, or
        future resale.
\end{itemize}
Valuations are independent private values: each bidder knows her own
type and valuation but not those of others.  Truthful bidding is
weakly dominant, so the auction price equals the second-highest
valuation among the~$n$ bidders.

%% ── 3.2  Primitives ─────────────────────────────────────────────

\subsection{Primitives}
\label{sec:model:primitives}

\begin{assumption}[Consumption Valuation]
\label{as:vc}
The consumption-buyer valuation is
\begin{equation}
  v_{c}(a^{*})
  = \alpha_{c}\, h(a^{*}),
  \qquad
  h(a^{*})
  = (a^{*})^{\gamma}\, e^{\,\gamma(1 - a^{*})},
  \quad \gamma > 0,
  \label{eq:vc}
\end{equation}
where $\alpha_{c} > 0$ is the maximum consumption value (attained at
$a^{*} = 1$).  The quality function~$h$ satisfies $h(0) = 0$,
$h(1) = 1$, is strictly increasing on $[0,1]$, strictly decreasing
on $(1,\infty)$, and $h(a^{*}) \to 0$ as $a^{*} \to \infty$.  The
parameter~$\gamma$ controls peakedness: higher~$\gamma$ implies a
sharper peak at maturity and faster post-maturity decay.
\end{assumption}

\begin{assumption}[Collector Valuation]
\label{as:vk}
The collector-buyer valuation is
\begin{equation}
  v_{k}(a^{*})
  = \bar{v}_{k} + \alpha_{k}\, a^{*},
  \qquad
  \bar{v}_{k} > 0,
  \quad \alpha_{k} \geq 0.
  \label{eq:vk}
\end{equation}
The parameter~$\bar{v}_{k}$ is the baseline collector value---the
``price floor'' reflecting scarcity, prestige, and storage of fine
wine---and $\alpha_{k}$ captures the incremental age premium.
\end{assumption}

\begin{assumption}[Buyer Composition]
\label{as:lambda}
The probability that a given bidder is type~$c$ is
\begin{equation}
  \lambda(a^{*})
  = \begin{cases}
      \bar{\lambda}
        & \text{if } a^{*} \leq 1, \\[4pt]
      \bar{\lambda}\, e^{-\delta(a^{*} - 1)}
        & \text{if } a^{*} > 1,
    \end{cases}
  \qquad
  \bar{\lambda} \in (0,1),
  \quad \delta > 0.
  \label{eq:lambda}
\end{equation}
Consumption buyers participate at a constant rate~$\bar{\lambda}$
before maturity and exit at rate~$\delta$ afterward, reflecting
declining interest in a wine that is past its drinking peak.
\end{assumption}

\noindent
The model has six parameters:
$(\alpha_{c},\, \bar{v}_{k},\, \alpha_{k},\, \gamma,\, \delta,\,
\bar{\lambda})$.  We suppress dependence on~$n$ (number of bidders)
in the notation; all results hold for any fixed $n \geq 2$.

%% ── 3.3  Equilibrium Price ──────────────────────────────────────

\subsection{Equilibrium Price}
\label{sec:model:equil}

Under truthful bidding, the auction price equals the second-highest
valuation.  Because each bidder's valuation is deterministic
conditional on type, the $n$ bids are i.i.d.\ draws from the
two-point distribution
$\{v_{c}(a^{*}),\, v_{k}(a^{*})\}$ with probabilities
$\{\lambda(a^{*}),\, 1 - \lambda(a^{*})\}$.

\medskip

\noindent\textbf{Definition (Equilibrium price).}
The expected auction price at normalised age~$a^{*}$ is
\begin{equation}
  P(a^{*})
  = \mathbb{E}\!\left[\,V_{(n-1:n)} \,\big|\, a^{*}\,\right],
  \label{eq:price}
\end{equation}
where $V_{(n-1:n)}$ denotes the second-order statistic of the~$n$
valuations.

\medskip

Write $\lambda = \lambda(a^{*})$, $v_{c} = v_{c}(a^{*})$,
$v_{k} = v_{k}(a^{*})$.  When $v_{c} \geq v_{k}$ (which holds near
maturity under Assumption~\ref{as:vc}--\ref{as:vk}), the
second-highest bid equals~$v_{c}$ if and only if at least two of the
$n$~bidders are type~$c$.  Letting
$q_{0} = (1-\lambda)^{n}$ and
$q_{1} = n\,\lambda\,(1-\lambda)^{n-1}$:
\begin{equation}
  P(a^{*})
  = v_{c}\!\left(1 - q_{0} - q_{1}\right)
  + v_{k}\!\left(q_{0} + q_{1}\right)
  \qquad \text{when } v_{c} \geq v_{k}.
  \label{eq:price_vc_dom}
\end{equation}
When $v_{k} > v_{c}$ (deep post-maturity), the symmetric expression
holds with $r_{0} = \lambda^{n}$ and
$r_{1} = n\,(1-\lambda)\,\lambda^{n-1}$:
\begin{equation}
  P(a^{*})
  = v_{k}\!\left(1 - r_{0} - r_{1}\right)
  + v_{c}\!\left(r_{0} + r_{1}\right)
  \qquad \text{when } v_{k} > v_{c}.
  \label{eq:price_vk_dom}
\end{equation}
In either case, $P(a^{*})$ is a weighted average of the two
valuations, with the weight on the higher valuation equal to the
probability that the second-order statistic equals that value.

%% ── 3.4  Main Results ───────────────────────────────────────────

\subsection{Main Results}
\label{sec:model:results}

\begin{proposition}[Non-Monotone Price-Age Profile]
\label{prop:nonmono}
Suppose\, $\alpha_{c} > \bar{v}_{k} + \alpha_{k}$
(consumption value at maturity exceeds collector value at maturity),
$\gamma > 0$, $\delta > 0$, and $n \geq 2$.  Then there exist
$a^{*}_{\mathrm{peak}}$ and $a^{*}_{\mathrm{trough}}$ with
$1 \leq a^{*}_{\mathrm{peak}} < a^{*}_{\mathrm{trough}}$ such that:
\begin{enumerate}
  \item[\emph{(i)}]
    $P(a^{*})$ is strictly increasing on
    $[0,\, a^{*}_{\mathrm{peak}}]$.
  \item[\emph{(ii)}]
    $P(a^{*})$ is strictly decreasing on
    $[a^{*}_{\mathrm{peak}},\, a^{*}_{\mathrm{trough}}]$.
  \item[\emph{(iii)}]
    $P(a^{*})$ is eventually strictly increasing for~$a^{*}$
    sufficiently large.
\end{enumerate}
\end{proposition}

\begin{proof}[Proof sketch]
Three forces shape the price path.
\emph{Pre-maturity} ($a^{*} \leq 1$): both $v_{c}$ and~$v_{k}$ are
increasing and $\lambda$ is constant, so $P$ is increasing.
\emph{Post-maturity transition} ($a^{*}$ slightly above~$1$):
$v_{c}$ decays at rate~$\gamma$ (Assumption~\ref{as:vc}) and
$\lambda$ decays at rate~$\delta$ (Assumption~\ref{as:lambda}).
Both effects reduce the probability that the second-highest bidder
is type~$c$, and since $v_{c} > v_{k}$ in a neighbourhood of
$a^{*} = 1$, the price falls---the \emph{composition effect}
dominates.
\emph{Recovery} ($a^{*} \to \infty$): $\lambda \to 0$ and
$v_{c} \to 0$, so the auction is populated entirely by collectors.
The price converges to~$v_{k}(a^{*})$, which is increasing.
The trough is the age at which the declining composition effect is
exactly offset by the rising collector premium.
A formal proof is in Appendix~\ref{app:proofs}.
\end{proof}

\medskip

The next result explains why the trough is absent for Burgundy Grand
Cru.

\begin{proposition}[Trough Disappearance]
\label{prop:notrough}
Define the collector-dominance ratio
$\rho \equiv \bar{v}_{k} / \alpha_{c}$.
There exists a threshold~$\rho^{*} \in (0,1)$ such that if
$\rho > \rho^{*}$, then $P(a^{*})$ is weakly increasing for all
$a^{*} \geq 0$.
\end{proposition}

\begin{proof}[Proof sketch]
When $\rho \geq 1$, we have $\bar{v}_{k} \geq \alpha_{c}$, so
$v_{k}(a^{*}) \geq \bar{v}_{k} \geq \alpha_{c} \geq v_{c}(a^{*})$
for all~$a^{*}$.  Collectors dominate every auction, and the price
tracks the increasing function~$v_{k}$.
For $\rho < 1$ but $\rho > \rho^{*}$, there exists a narrow interval
near $a^{*} = 1$ where $v_{c}$ exceeds~$v_{k}$, but the interval is
too short and the consumption share~$\bar{\lambda}$ too small for
the composition effect to generate a net price decline.
The threshold~$\rho^{*}$ depends on $(\gamma, \delta, \bar{\lambda},
\alpha_{k}, n)$; a formal characterisation is in
Appendix~\ref{app:proofs}.
\end{proof}

\noindent
The ratio~$\rho$ has a direct empirical counterpart: the share of
young auction lots (normalised age below~0.4) whose per-bottle price
exceeds~\$200.  This share proxies for the intensity of collector
demand at early ages.  Table~\ref{tab:rho_mapping} summarises the
cross-sectional mapping.

\begin{table}[htbp]
\centering
\caption{Collector-Dominance Ratio and Trough Presence}
\label{tab:rho_mapping}
\begin{tabular}{lccc}
\toprule
Category & Young lots $>$ \$200 & Implied $\rho$ & Trough \\
\midrule
Burgundy Grand Cru     & 51.5\% & High ($> \rho^{*}$) & Absent \\
Bordeaux Crus Class\'es & 39.0\% & Moderate ($< \rho^{*}$) & $-26.2\%$ \\
Burgundy Premier Cru   & 11.2\% & Low ($\ll \rho^{*}$) & $-22.4\%$ \\
\bottomrule
\end{tabular}
\end{table}

\begin{proposition}[Comparative Statics on Trough Depth]
\label{prop:compstat}
For a wine category with $\rho < \rho^{*}$, define the trough depth
$\Delta \equiv P(a^{*}_{\mathrm{peak}}) -
P(a^{*}_{\mathrm{trough}})$.  Then:
\begin{enumerate}
  \item[\emph{(i)}]
    $\partial \Delta / \partial \bar{\lambda} > 0$: a higher
    consumption-buyer share deepens the trough.
  \item[\emph{(ii)}]
    $\partial \Delta / \partial \bar{v}_{k} < 0$: a higher collector
    floor compresses the trough.
  \item[\emph{(iii)}]
    $\partial \Delta / \partial \gamma > 0$: a sharper consumption
    peak deepens the trough.
  \item[\emph{(iv)}]
    $\partial \Delta / \partial \delta > 0$: faster consumption-buyer
    exit deepens the trough.
\end{enumerate}
\end{proposition}

\begin{proof}[Proof sketch]
Part~\emph{(i)}: higher~$\bar{\lambda}$ increases competition among
consumption buyers near maturity, raising
$P(a^{*}_{\mathrm{peak}})$.  After maturity, the same buyers exit,
so $P(a^{*}_{\mathrm{trough}})$ falls---the gap widens.
Part~\emph{(ii)}: higher~$\bar{v}_{k}$ raises the price floor set
by collectors, truncating the price decline at the trough.
Parts~\emph{(iii)}--\emph{(iv)} accelerate the post-maturity
collapse of consumption demand, steepening the transition from peak
to trough.
Formal proofs are in Appendix~\ref{app:proofs}.
\end{proof}

%% ── 3.5  Empirical Predictions ──────────────────────────────────

\subsection{Empirical Predictions}
\label{sec:model:predictions}

The model yields four testable implications that we take to the data
in Sections~\ref{sec:empirics}--\ref{sec:results}.

\medskip

\noindent\textbf{Prediction~1} (Non-monotone profile).
For wine categories in which consumption buyers constitute a
substantial share of bidders ($\rho < \rho^{*}$), the log price-age
profile should rise pre-maturity, fall post-maturity, and recover at
antique ages.

\medskip

\noindent\textbf{Prediction~2} (Trough disappearance).
Wine categories where the collector-dominance ratio is sufficiently
high---observable as a high share of young lots priced above \$200
per bottle---should exhibit a monotonically increasing price-age
profile with no post-maturity trough.

\medskip

\noindent\textbf{Prediction~3} (Monotone pre-maturity rise).
All categories should show rising prices pre-maturity, because both
consumption and collector valuations are increasing before $a^{*}=1$.

\medskip

\noindent\textbf{Prediction~4} (Cross-sectional trough depth).
Across categories that exhibit a trough, its depth should be
increasing in the consumption-buyer share and decreasing in the
collector-value floor.

%% ── 3.6  Discussion ─────────────────────────────────────────────

\subsection{Discussion}
\label{sec:model:discussion}

\paragraph{Relationship to \citet{lovo2018}.}
Our model shares the two-type buyer structure with the art-market
model of \citet{lovo2018}, but differs in three key respects.
First, wine has a finite consumption life governed by the maturity
date~$\tau$, which generates the single-peaked consumption valuation
$v_{c}(a^{*})$ central to the non-monotone price profile.
Art has no biological expiry, so the consumption-type valuation in
\citet{lovo2018} is constant.
Second, the maturity date varies across wines and is (imperfectly)
observable via expert tasting notes, providing cross-sectional
variation in~$a^{*}$ that identifies the model.
Third, we adopt a static auction mechanism appropriate for the
institutional setting, rather than the dynamic search-and-matching
framework of \citet{lovo2018}.

\paragraph{Why arbitrageurs do not eliminate the trough.}
If collectors could costlessly buy at the trough and resell at
antique ages, forward-looking bidding would compress the price dip.
The persistence of the trough in the data points to significant
frictions: storage costs (approximately 1--2\,\% of value per year
for temperature-controlled storage), authentication risk (provenance
uncertainty rises with the number of transactions), and thin markets
(low trading frequency for post-maturity wines reduces liquidity).

\paragraph{Endogenous supply.}
The baseline model takes supply as given.  If sellers anticipate
lower prices in the trough zone, they may delay consignment,
reducing supply and partially offsetting the price decline.  The fact
that the trough persists suggests that seller delay, while plausible,
does not fully eliminate the composition-driven dip.

\paragraph{Structural estimation.}
The six model parameters
$(\alpha_{c}, \bar{v}_{k}, \alpha_{k}, \gamma, \delta,
\bar{\lambda})$ can be estimated by matching the observed price-age
profiles across wine categories using minimum distance, with 7
moment conditions (8 age bins, young as baseline) and the full
HC1-robust covariance matrix as optimal weight.
Appendix~\ref{app:structural} describes the estimation procedure and
will report these estimates for all three grand cru categories, together
with a $J$-test of overidentification and parametric bootstrap
confidence sets.
\textit{[Placeholder: structural estimation results to be completed.
The estimated $\rho = \bar{v}_{k}/\alpha_{c}$ is expected to exceed
unity for Burgundy Grand Cru, confirming Proposition~\ref{prop:notrough}.
]}